\section{Winner-Take-All Global-Trace Holonomy SDP}
\label{sec:winner-take-all}

This section gives an explicit SDP family with bounded scalar incidence to
which the condensation theory of \cref{sec:condensation} applies.  It couples
$G$ groups, each a copy of the qutrit-holonomy component of
\cref{sec:sdp-intrinsic}.  Each group has \(N=n_{\rm bit}\) private signs,
and its holonomy is their parity.  A single trace constraint shared by all
groups makes the optimal value the least minimum eigenvalue of the group
costs, and a chain of free scalar accumulators keeps every row and column of
the equality matrix sparse.

\begin{center}
\begin{tabular}{@{}ll@{}}
\toprule
quantity & conclusion \\ \midrule
value queries & $Q=\Theta(N\sqrt G)$, $R=\Theta(NG)$ for error $<1/20$ \\
one-odd central path & $p=\Theta(1),\Theta(G^{-1/2}),\Theta(G^{-1})$ \\
 & below, at, and above $\alpha_*=2/45$ \\
reduced-Hessian conditioning & $\Theta(G),\Theta(G),\Theta(1)$ in the same phases \\
coherent central-state preparation & $\Theta(N\sqrt G),\Theta(NG^{1/4}),O(N)$ \\
end-to-end algorithm & search for the winner, then solve only its group \\
\bottomrule
\end{tabular}
\end{center}
Here $Q$ and $R$ are the quantum and randomized query complexities of the
optimal value, $p$ is the central-path trace mass of the odd group when
exactly one group is odd, and the phases refer to $\alpha=\tau G$
(\cref{sec:wta-central-path}).  The preparation costs assume exactly one
winner.  At the public starting point the root Hessian has condition number
one, but the Newton-direction state does not reveal the winner
(\cref{rem:wta-dilution}).

\subsection{A bounded-incidence product SDP}
\label{sec:wta-construction}

Set
\begin{equation}
 K=16N,\qquad P=2K+N=33N,
 \qquad u=\frac1{10},\qquad q=\frac KP=\frac{16}{33},
 \qquad \gamma=\frac uq=\frac{33}{160}.
 \label{eq:wta-parameters}
\end{equation}
For $a<b$, write $H_{ab}=e_ae_b^\top+e_be_a^\top$, and put
$D_s=\diag(s,1,1)$ for $s\in\{-1,+1\}$.  Group $g$ has private signs
$\sigma_{g,1},\ldots,\sigma_{g,N}$ and holonomy
\begin{equation}
 h_g=\prod_{j=1}^N\sigma_{g,j}.
 \label{eq:wta-holonomy}
\end{equation}
Group $g$ has $P$ blocks $X_{g,0},\ldots,X_{g,P-1}\in\Splus^3$, joined in a
path.  For the edge from $i-1$ to $i$, let $G_{g,i}=I$ except that
$G_{g,K+j-1}=D_{\sigma_{g,j}}$ for $1\le j\le N$, and impose the copy
equations, each as six isometric-svec scalar equations:
\begin{equation}
 X_{g,i}=G_{g,i}X_{g,i-1}G_{g,i}^\top
 \qquad(1\le i<P).
 \label{eq:wta-copy-rows}
\end{equation}
If $R_{g,0}=I$ and $R_{g,i}=G_{g,i}\cdots G_{g,1}$, these equations state that
\begin{equation}
 X_{g,i}=R_{g,i}Z_gR_{g,i}^\top,\qquad Z_g=X_{g,0}\succeq0.
 \label{eq:wta-root-lift}
\end{equation}
The first $K$ matrices $R_{g,i}$ equal $I$, and the last $K$ equal
$D_{h_g}$.

The public block costs are
\begin{equation}
 C_{g,i}=\begin{cases}
 3I+uH_{23}+\gamma H_{12},&0\le i<K,\\
 3I+uH_{23},&K\le i<K+N,\\
 3I+uH_{23}+\gamma H_{13},&K+N\le i<P.
 \end{cases}
 \label{eq:wta-local-costs}
\end{equation}
Because $\sqrt{u^2+\gamma^2}<1/4$, every $C_{g,i}$ has spectrum in
$(11/4,13/4)$.  Introduce free public scalars $t_1,\ldots,t_G$ and add
\begin{align}
 t_1-\tr X_{1,0}&=0,\nonumber\\
 t_g-t_{g-1}-\tr X_{g,0}&=0 &&(2\le g\le G),\nonumber\\
 t_G&=1. \label{eq:wta-accumulator}
\end{align}
For $G=1$, the first and last rows impose $t_1=\tr X_{1,0}=1$.
The SDP minimizes
\begin{equation}
 f(X)=\frac1P\sum_{g=1}^G\sum_{i=0}^{P-1}
                  \inner{C_{g,i}}{X_{g,i}}
 \label{eq:wta-objective}
\end{equation}
subject to \eqref{eq:wta-copy-rows}--\eqref{eq:wta-accumulator} and
$X_{g,i}\succeq0$.

There are
\begin{equation}
 L=GP=33GN
 \label{eq:wta-block-count}
\end{equation}
blocks in $\Splus^3$.  A copy row has two scalar nonzeros.  The first
accumulator row has four, an interior accumulator row has five, and the
final row $t_G=1$ has one.
Every matrix coordinate and every free-scalar column occurs in at most two
rows.  Thus the equality matrix has scalar row sparsity at most five and
scalar column sparsity at most two.  Its support, right-hand side, and all
costs are public; only the $12$ and $13$ copy coefficients on the $N$ signed
edges of each path depend on the input.

Eliminating the path and accumulator variables leaves the roots $Z_g$,
whose feasible set is exactly
\begin{equation}
 \mathcal B_G=\left\{(Z_1,\ldots,Z_G):Z_g\succeq0,
                       \ \sum_{g=1}^G\tr Z_g=1\right\}.
 \label{eq:wta-root-feasible-set}
\end{equation}
The public point
\begin{equation}
 X_{g,i}^0=\frac1{3G}I,\qquad t_g^0=\frac gG
 \label{eq:wta-public-point}
\end{equation}
is strictly primal feasible.  With all equality multipliers zero, the dual
PSD slacks are $C_{g,i}/P\succ0$ and stationarity in every free scalar is
zero.  Hence the formulation is strictly primal--dual feasible in the sense
of \cref{def:strict-feasibility}.

\subsection{Exact spectral minimum}
\label{sec:wta-optimum}

In terms of the roots, the objective \eqref{eq:wta-objective} becomes
\begin{equation}
 \overline C_h=3I+u(H_{12}+H_{23}+hH_{13}),
 \qquad
 \min_{Z\in\mathcal B_G}\sum_g\inner{\overline C_{h_g}}{Z_g}.
 \label{eq:wta-effective-cost}
\end{equation}
Indeed, the first and last $K$ blocks of each path contribute
$K\gamma/P=u$ to $H_{12}$ and $H_{13}$, respectively, and
conjugation by $D_h$ changes $H_{13}$ to $hH_{13}$.  The two effective costs
have spectra
\begin{equation}
 \operatorname{spec}(\overline C_+)
   =\left\{\frac{16}5,\frac{29}{10},\frac{29}{10}\right\},
 \qquad
 \operatorname{spec}(\overline C_-)
   =\left\{\frac{31}{10},\frac{31}{10},\frac{14}5\right\}.
 \label{eq:wta-cost-spectra}
\end{equation}
These follow either by diagonalizing the two signed triangle adjacency
matrices, whose eigenvalues are $(2,-1,-1)$ and $(1,1,-2)$, or from the
intrinsic component analysis in \cref{sec:sdp-intrinsic}.

For any feasible roots, Rayleigh--Ritz gives
\begin{equation}
 \sum_g\inner{\overline C_{h_g}}{Z_g}
 \ge\sum_g\lambda_{\min}(\overline C_{h_g})\tr Z_g
 \ge\min_g\lambda_{\min}(\overline C_{h_g}).
 \label{eq:wta-rayleigh-ritz}
\end{equation}
Equality holds when all trace is placed on a minimum eigenvector of a
minimizing group.  Therefore
\begin{equation}
 \boxed{\operatorname{OPT}=
 \begin{cases}
 29/10,&h_g=+1\text{ for every }g,\\
 14/5,&h_g=-1\text{ for at least one }g.
 \end{cases}}
 \label{eq:wta-optimum}
\end{equation}
The gap is $1/10$, so the optimal value encodes
$\operatorname{OR}_G\circ\operatorname{PARITY}_N$: whether some $h_g=-1$.

The root feasible set \eqref{eq:wta-root-feasible-set} has no lift over a
finite product of second-order cones.  Its conic hull is $(\Splus^3)^G$.  If
\eqref{eq:wta-root-feasible-set} had such a lift, homogenizing the compact lift and
projecting to one factor would give such a lift of $\Splus^3$, contrary to
Fawzi's theorem \cite{fawzi2018-on-representing-the-positive-semidefinite}.
The path-and-accumulator formulation projects linearly onto the root set, so
it has no finite SOC lift either.

\subsection{Sharp value-query complexity}
\label{sec:wta-query-complexity}

We use the raw sparse coefficient access of \cref{def:raw-sparse-access}.
In the canonical oracle, which returns isometric-svec coefficients at fixed
positions, a public reversible address map sends a requested
input-dependent position to $(g,j)$ and a public position to a dummy zero.
The value oracle XORs the sign of one of the two hidden coefficients
$\pm\sigma_{g,j}$ directly into its target register.  Thus one raw-sign query, which returns one $\sigma_{g,j}$, simulates
one coefficient query, coherently and also under inverse
or external control.  Conversely, one fixed input-dependent copy coefficient
reveals $\sigma_{g,j}$.  The two models
therefore have the same query counts.  A clean value oracle can
instead be used as a phase oracle by compute--phase--uncompute, at two calls
per query.

We instantiate \cref{thm:cond-unique-or} for parity blocks.
Let $E,O\subseteq\{0,1\}^N$ be the even- and odd-parity strings and let
\begin{equation}
 B_{xy}=\mathbf1[|x-y|_1=1]\qquad(x\in E,\ y\in O).
 \label{eq:wta-hypercube}
\end{equation}
The constant vectors show that $\norm B=N$, while each
$B\circ\Delta_j$ is a permutation matrix and has norm one.  Hence the inner
adversary value in \eqref{eq:cond-inner-adversary} is exactly $N$.
Applying \eqref{eq:cond-unique-or-adv} on the all-even versus
exactly-one-odd promise gives
\begin{equation}
 Q_{1/3}=\Omega(N\sqrt G).
 \label{eq:wta-quantum-lower}
\end{equation}
The hypercube matrix \eqref{eq:wta-hypercube} is thus an explicit
adversary witness for \cref{sec:condensation-query-lower}.

If $|\widetilde v-\operatorname{OPT}|\le\epsilon<1/20$, thresholding at
$57/20$ distinguishes the two cases in \eqref{eq:wta-optimum}.  The same
threshold works for target-relative error $1/100$, since
\begin{equation}
 \frac{101}{100}\frac{14}5<\frac{57}{20}
 <\frac{99}{100}\frac{29}{10}.
 \label{eq:wta-relative-threshold}
\end{equation}
At $\epsilon=1/20$, the midpoint $57/20$ is a valid answer for both values
and uses no query, so the additive threshold is sharp.

For the randomized lower bound, let $\mathcal D_0$ make all groups
independent and uniform conditional on even parity; let $\mathcal D_1$
choose a uniform group $J$, make it uniform conditional on odd parity, and
leave the others as in $\mathcal D_0$.  Couple adaptive decision trees by returning common fair answers to the first $N-1$ distinct positions in each
group and using the final unseen sign to enforce the promised parity.  The
transcripts can differ only if the algorithm queries all $N$ signs of
group $J$.  A depth-$r$ tree queries all signs of at most $r/N$ groups, so
\begin{equation}
 \operatorname{TV}(\mathsf{Tr}_{\mathcal D_0},
                    \mathsf{Tr}_{\mathcal D_1})\le\frac r{NG}.
 \label{eq:wta-yao-coupling}
\end{equation}
Yao's principle therefore gives $r\ge NG/3$ for worst-case success $2/3$;
reading every sign is a matching upper bound.  Consequently, for every
$\epsilon<1/20$,
\begin{equation}
 \boxed{Q_\epsilon(\operatorname{OPT})=\Theta(N\sqrt G),
 \qquad R_\epsilon(\operatorname{OPT})=\Theta(NG).}
 \label{eq:wta-value-complexity}
\end{equation}

\subsection{The Newton step at the public point}
\label{sec:wta-public-start}

Use the log-det barrier on the PSD blocks, but no barrier on the free
accumulators:
\begin{equation}
 \Phi_\mu(X)=f(X)-\mu\sum_{g,i}\log\det X_{g,i},
 \qquad \mu=\frac1{GP}.
 \label{eq:wta-start-barrier}
\end{equation}
At \eqref{eq:wta-public-point}, the PSD-block gradient and Hessian in
isometric-svec coordinates are
\begin{equation}
 g_{g,i}=\frac1P(C_{g,i}-3I),
 \qquad \nabla_X^2\Phi_\mu(X^0)=\frac{9G}{P}I.
 \label{eq:wta-start-gradient-hessian}
\end{equation}
These and all equality residuals are public; the input enters the unreduced
Newton system only through the signed copy constraints in its matrix.

After exact elimination of the copies and accumulators, a tangent
direction in root coordinates is $Y=(Y_1,\ldots,Y_G)$ with
$\sum_g\tr Y_g=0$.  The copy maps are orthogonal congruences and each root
has $P$ path copies, so
\begin{equation}
 \boxed{\nabla^2_{\rm root}\Phi_\mu(X^0)=9G I_{6G-1}.}
 \label{eq:wta-root-hessian}
\end{equation}
This root-coordinate Hessian has condition number one, including in
directions that move trace between groups.  With
\begin{equation}
 A_h=H_{12}+H_{23}+hH_{13},
 \label{eq:wta-Ah}
\end{equation}
the pulled-back gradient is $uA_{h_g}$, which is trace zero, and the exact
Newton directions are
\begin{equation}
 \Delta Z_g=-\frac{u}{9G}A_{h_g},\qquad \Delta t_g=0.
 \label{eq:wta-newton-direction}
\end{equation}
Since $\norm{A_h}_F^2=6$, the decrement is
\begin{equation}
 \lambda^2=\sum_g\inner{uA_{h_g}}{(9G)^{-1}uA_{h_g}}
 =\frac1{150},\qquad \lambda=\frac1{\sqrt{150}}.
 \label{eq:wta-decrement}
\end{equation}
The full step remains strictly feasible.  The eigenvalues inside the
parentheses below are $14/45,31/90,31/90$ for $h=+1$ and
$29/90,29/90,16/45$ for $h=-1$.  Consequently,
\begin{equation}
 Z_g^0+\Delta Z_g=\frac1G\left(\frac13I-\frac u9A_{h_g}\right)
 \succeq\frac{14}{45G}I.
 \label{eq:wta-full-step-margin}
\end{equation}
On the root slice the barrier weight is $\tau=\mu P=1/G$ (see
\eqref{eq:wta-root-center-problem}), which is below one for every $G\ge2$;
then $\Phi_\mu$ is not standard self-concordant there, and the usual
full-step decrement estimate does not apply to $\lambda$.  We instead bound the next
decrement directly.  With
$B_h=\tfrac13I-\tfrac u9A_h$, the full step is $\widehat Z_g=B_{h_g}/G$ and
its slice gradient is $\overline C_{h_g}-B_{h_g}^{-1}$.  On the trace slice
the Newton decrement is the minimum over the trace multiplier of the
Hessian-inverse norm of the shifted gradient, so taking that multiplier to
be zero gives an upper bound:
\begin{equation}
 \lambda_{\rm next}^2
 \le\frac1G\sum_g\tr\bigl[(B_{h_g}\overline C_{h_g}-I)^2\bigr]
 =\frac{u^4}{81}\tr A_h^4=\frac{2u^4}{9},
 \qquad
 \lambda_{\rm next}\le\frac{\sqrt2}{300},
 \label{eq:wta-next-decrement}
\end{equation}
because $\overline C_h=3I+uA_h$ gives
$B_h\overline C_h-I=-\tfrac{u^2}9A_h^2$ and $\tr A_h^4=18$ for both
parities.  This bound is uniform in $G$ and in the parity pattern.

\begin{remark}[The Newton-direction state hides the winner]
\label{rem:wta-dilution}
The normalized root-direction state is
\begin{equation}
 \ket{\Delta(h)}=\frac1{\sqrt G}\sum_{g=1}^G\ket g\ket{a_{h_g}},
 \qquad
 \ket{a_h}=\frac{\ket{\operatorname{svec}(A_h)}}{\sqrt6}.
 \label{eq:wta-newton-state}
\end{equation}
Looping once over the $N$ raw signs and applying phase kickback only to the
$13$ coordinate prepares this state exactly with $N$ queries.  Conversely,
the density-matrix entry between group~1's $12$ and $13$ coordinates is
$h_1/(3G)$.  Since each density-matrix entry of a $q$-query
output has degree at most $2q$ in the input signs, exact preparation needs $q\ge N/2$; the same holds at trace error
$\delta<1/(6G)$.  Hence exact preparation costs $\Theta(N)$.

Moreover, the all-even and one-odd states have overlap $1-2/(3G)$ and trace
distance
\begin{equation}
 \frac{2\sqrt{3G-1}}{3G}=\Theta(G^{-1/2}).
 \label{eq:wta-newton-distance}
\end{equation}
Therefore a constant-error Newton state cannot decide the unique-winner
promise with constant bias.  The $\Omega(N\sqrt G)$ lower bound applies to
the value and to optimizer outputs that reveal the winner, not to this
state.
\end{remark}

The condition number one in \eqref{eq:wta-root-hessian} holds in root
coordinates only.  A feasible root direction lifts to accumulator changes
equal to prefix sums of traces, which is not an isometry, and the free
scalars have zero barrier Hessian.  We therefore do not claim condition
number one for the ambient augmented Hessian or the saddle-point KKT
matrix.

\subsection{Approximate optimizers}
\label{sec:wta-approximate-optimizers}

For exactly feasible roots on an instance with at least one odd group, put
$m_-=\sum_{g:h_g=-1}\tr Z_g$.  If the objective is at most
$14/5+\delta_{\rm obj}$, \eqref{eq:wta-cost-spectra} and
\eqref{eq:wta-rayleigh-ritz} give
\begin{equation}
 \frac{14}5+\frac{1-m_-}{10}\le \sum_g
 \inner{\overline C_{h_g}}{Z_g}
 \le\frac{14}5+\delta_{\rm obj},
 \qquad
 \boxed{m_-\ge1-10\delta_{\rm obj}.}
 \label{eq:wta-exact-optimizer-mass}
\end{equation}
When the instance contains both an odd and an even group and
$0\le\delta_{\rm obj}\le1/10$, the constant is sharp: put trace $1-10\delta_{\rm obj}$ in an odd
and $10\delta_{\rm obj}$ in an even minimum eigenspace.

A similar bound holds for approximate solutions of the sparse formulation,
which may violate its equalities.  Define the copy residuals
\begin{equation}
 E_{g,i}=X_{g,i}-G_{g,i}X_{g,i-1}G_{g,i}^\top
 \label{eq:wta-copy-residual}
\end{equation}
and let $r_1,\ldots,r_G,r_{\rm top}$ be the residuals in
\eqref{eq:wta-accumulator}.  Suppose $X_{g,i}\succeq0$ and
\begin{equation}
 \sum_{g,i\ge1}\norm{E_{g,i}}_F^2+
 \sum_{g=1}^G r_g^2+r_{\rm top}^2\le\epsilon^2,
 \qquad f(X)\le\operatorname{OPT}+\delta_{\rm obj}.
 \label{eq:wta-approx-contract}
\end{equation}
Put $Z_{g,i}=R_{g,i}^\top X_{g,i}R_{g,i}$ (each block in its root frame) and
telescope along the paths and the accumulator chain.
Cauchy--Schwarz and the spectral bound on $C_{g,i}$ give
\begin{align}
 \left|\sum_g\tr Z_{g,0}-1\right|&\le\sqrt{G+1}\,\epsilon,
 \label{eq:wta-root-trace-error}\\
 \left|\frac1P\sum_{g,i}\tr X_{g,i}-1\right|&\le\sqrt{GP+1}\,\epsilon,
 \label{eq:wta-physical-trace-error}\\
 \left|f(X)-\sum_g\inner{\overline C_{h_g}}{Z_{g,0}}\right|
 &<\frac{13}{4}\sqrt{GP}\,\epsilon.
 \label{eq:wta-objective-lift-error}
\end{align}
Combining these estimates with \eqref{eq:wta-cost-spectra} gives, for the
even root mass $m_+$ on a yes-instance,
\begin{equation}
 m_+\le10\delta_{\rm obj}+28\sqrt{G+1}\,\epsilon
              +\frac{65}{2}\sqrt{GP}\,\epsilon.
 \label{eq:wta-robust-even-mass}
\end{equation}

\begin{theorem}[Robust winner recovery]
\label{thm:wta-robust-optimizer}
Suppose an algorithm prepares a state within trace distance $1/100$ of the
trace-normalized physical density
\begin{equation}
 \rho(X)=\frac{\bigoplus_{g,i}X_{g,i}}
                 {\sum_{g,i}\tr X_{g,i}}
 \label{eq:wta-optimizer-density}
\end{equation}
for PSD blocks and free accumulators satisfying
\eqref{eq:wta-approx-contract}, where
\begin{equation}
 \delta_{\rm obj}\le\frac1{100},
 \qquad \epsilon\le\frac1{1000\sqrt{GP}}.
 \label{eq:wta-optimizer-tolerances}
\end{equation}
Then a fixed input-independent measurement decides whether any
group is odd with constant bias.  On every yes-instance, measuring the group register returns
an actual odd group with probability greater than $5/6$.  Consequently any
such first preparation (in the sense of \cref{def:amortized-accounting})
uses $\Omega(N\sqrt G)$ raw or canonical coefficient
queries.
\end{theorem}

\begin{proof}
Set $s=\sqrt{GP}\epsilon\le1/1000$.  Since $P\ge33$,
$\sqrt{G+1}\epsilon<s/4$.  Equations
\eqref{eq:wta-physical-trace-error} and
\eqref{eq:wta-robust-even-mass} put the physical trace of the even groups
below
$1/10+(81/2)s$ and the total trace above $1-(33/32)s$.  Their ratio is less
than $1/7$; trace error $1/100$ leaves odd-group probability above $5/6$.
For existence alone, use the input-independent observable
\begin{equation}
 M=\frac1{2/5}\left(\bigoplus_{g,i}C_{g,i}-\frac{57}{20}I\right)
 \label{eq:wta-existence-observable}
\end{equation}
which has norm at most one.  The same three telescoping estimates put
$f(X)/\overline T<2.82$ on yes-instances and $>2.89$ on no-instances,
where $\overline T=P^{-1}\sum_{g,i}\tr X_{g,i}$.  Equivalently,
$\tr(M\rho(X))<-3/40$ and $>1/10$, respectively.
The POVM $(I\pm M)/2$, repeated a constant number of times, therefore
decides the promise.  Now apply \eqref{eq:wta-quantum-lower}.
\end{proof}

The order $1/\sqrt{GP}$ of the residual tolerance in
\eqref{eq:wta-optimizer-tolerances} is necessary.  On an all-even base input,
flip each of the $GN$ raw signs in turn, take the neighboring instance's
exact odd optimizer, and average the resulting physical solutions.  This
PSD mixture has objective $14/5$, exact accumulator rows, and only the $GN$
flipped-edge copy residuals.  With
$r_-=(1,-1,1)^\top/\sqrt3$ and $Z_-^*=r_-r_-^\top$, each nonzero residual before
averaging has Frobenius norm
$\norm{D_-Z_-^*D_--Z_-^*}_F=4/3$.  The total Euclidean residual is therefore
\begin{equation}
 \frac4{3\sqrt{GN}}=\frac{4\sqrt{33}}{3\sqrt{GP}}.
 \label{eq:wta-residual-sharpness}
\end{equation}
Yet on the all-even input its density makes
\eqref{eq:wta-existence-observable} report a winner.  Feasibility alone is
also insufficient: the public point \eqref{eq:wta-public-point} is exactly
feasible for every input.  Hence both conditions in
\cref{thm:wta-robust-optimizer} are needed.

\subsection{Exact central path and condensation}
\label{sec:wta-central-path}

Eliminating a length-$P$ path multiplies the physical barrier coefficient by
$P$.  Thus, with $\tau=\mu P$, the root problem is
\begin{equation}
 \min_{Z_g\succ0,\ \sum_g\tr Z_g=1}
 \left\{\sum_g\inner{\overline C_{h_g}}{Z_g}
                    -\tau\sum_g\log\det Z_g\right\}.
 \label{eq:wta-root-center-problem}
\end{equation}
Its exact center is
\begin{equation}
 Z_g(\tau)=\tau(\overline C_{h_g}-\lambda I)^{-1},
 \qquad
 1=\tau\sum_g\tr(\overline C_{h_g}-\lambda I)^{-1},
 \label{eq:wta-exact-center}
\end{equation}
with the unique $\lambda<\min_g\lambda_{\min}(\overline C_{h_g})$.
This instantiates \cref{prop:cond-exact-center}.

On the all-even input, every group has trace exactly $1/G$.  With
$y=29/10-\lambda$ and $\alpha=\tau G$, the scalar equation and its positive
solution are
\begin{align}
 1&=\alpha\left(\frac2y+\frac1{y+3/10}\right),
 \label{eq:wta-even-secular}\\
 y_\alpha&=\frac{3\alpha-3/10+
 \sqrt{(3/10-3\alpha)^2+(12/5)\alpha}}2.
 \label{eq:wta-even-root}
\end{align}

Now suppose group~1 is the unique odd group and put $x=14/5-\lambda$.  The
winner and loser resolvents are
\begin{equation}
 A(x)=\frac1x+\frac2{x+3/10},
 \qquad
 B(x)=\frac1{x+2/5}+\frac2{x+1/10}.
 \label{eq:wta-resolvents}
\end{equation}
The exact central path is determined by
\begin{equation}
 \boxed{1=\tau\{A(x)+(G-1)B(x)\}},
 \qquad
 \boxed{p_G=\tau A(x),\quad1-p_G=\tau(G-1)B(x)}.
 \label{eq:wta-one-odd-center}
\end{equation}
Direct algebra gives
\begin{equation}
 A(x)-B(x)=
 \frac{3/250}{x(x+1/10)(x+3/10)(x+2/5)}>0
 \qquad(x>0).
 \label{eq:wta-resolvent-order}
\end{equation}
Together with \eqref{eq:wta-one-odd-center}, this yields
$p_G/(1-p_G)>1/(G-1)$, and hence $p_G\ge1/G$.  This implication is special
to the present spectra and is not part of the general two-cost model of
\cref{sec:condensation-model}.
Since $B$ is decreasing on $x>0$ and $B(0)=45/2$,
\eqref{eq:wta-one-odd-center} also gives
\begin{equation}
 p_G\ge1-\frac{45}{2}\tau(G-1).
 \label{eq:wta-finite-mass}
\end{equation}

The constants of \cref{thm:condensation-threshold} are
\begin{equation}
 B_0=B(0)=\frac{45}{2},\qquad
 \beta=-B'(0)=\frac{825}{4},\qquad
 \boxed{\alpha_*=B_0^{-1}=\frac2{45}}.
 \label{eq:wta-critical-constants}
\end{equation}
Applying \cref{thm:condensation-threshold} gives, for fixed $\alpha>0$ and
$G\to\infty$,
\begin{align}
 0<\alpha<2/45:&\quad
 p_G\longrightarrow1-\frac{45}{2}\alpha,
 &\frac{x}{\tau}&\longrightarrow
       \frac1{1-(45/2)\alpha},
 \label{eq:wta-subcritical}\\
 \alpha=2/45:&\quad
 \sqrt G\,p_G\longrightarrow\frac{\sqrt{33}}9,
 &\sqrt G\,x&\longrightarrow\frac2{\sqrt{825}},
 \label{eq:wta-critical}\\
 \alpha>2/45:&\quad
 Gp_G\longrightarrow\alpha A(x_\alpha),
 &B(x_\alpha)&=\frac1\alpha.
 \label{eq:wta-supercritical}
\end{align}
In the supercritical case, \eqref{eq:wta-resolvent-order} implies the strict
coefficient inequality $\alpha A(x_\alpha)>\alpha B(x_\alpha)=1$.
The coefficient $\sqrt{33}/9$ follows from
$\alpha_*\sqrt\beta=(2/45)(\sqrt{825}/2)$.

At the center, the reduced Hessian in the Frobenius inner product is
\begin{equation}
 \mathcal H_\tau[Y]_g=\frac1\tau D_gY_gD_g,
 \qquad D_g=\overline C_{h_g}-\lambda I,
 \qquad \sum_g\tr Y_g=0.
 \label{eq:wta-central-hessian}
\end{equation}
For one odd group, the spectral gaps are
$(x,x+3/10,x+3/10)$ in the winner and
$(x+1/10,x+1/10,x+2/5)$ in every loser.  For $G\ge3$,
\cref{lem:cond-hessian-sandwich} specializes to
\begin{equation}
 \kappa_{\rm red}\asymp
 \frac{(x+2/5)^2}{
  \min\{x(x+3/10),\ x^2+(x+1/10)^2/G\}}.
 \label{eq:wta-kappa-finite}
\end{equation}
Therefore \cref{thm:cond-conditioning-phases} gives
\begin{equation}
 \kappa_{\rm red}=\begin{cases}
 \Theta(G),&0<\alpha\le2/45,\\
 \Theta(1),&\alpha>2/45.
 \end{cases}
 \label{eq:wta-conditioning-transition}
\end{equation}
On the all-even input and for $G\ge2$, by contrast,
\begin{equation}
 \kappa_{\rm even}=\left(\frac{y_\alpha+3/10}{y_\alpha}\right)^2
 =\Theta_\alpha(1).
 \label{eq:wta-even-conditioning}
\end{equation}
The bound \eqref{eq:cond-mc3} simplifies here because $\min\{1,(2/5)/(3/10)\}=1$: for $G\ge3$,
\begin{equation}
 \kappa_{\rm red}\ge\frac{(G-1)p_G}{1-p_G}.
 \label{eq:wta-mass-conditioning}
\end{equation}
Thus a constant winner mass forces $\kappa_{\rm red}=\Omega(G)$ without
preconditioning.  This concerns the central path only and gives no lower
bound for methods with an input-dependent preconditioner.  The corresponding
visibility bounds are \cref{eq:cond-cv2,eq:cond-cv4}.

For approximate centers, we measure centrality by the dimensionless local
residual of \cref{thm:cond-robust}, not an unscaled stationarity error.
If positive definite
$\widetilde Z_g$ and a multiplier $\widetilde\lambda<14/5$ satisfy
\begin{equation}
 \left\|\tau^{-1}\widetilde Z_g^{1/2}
 (\overline C_{h_g}-\widetilde\lambda I)
 \widetilde Z_g^{1/2}-I\right\|_{\rm op}\le\eta<1,
 \qquad
 \left|\sum_g\tr\widetilde Z_g-1\right|\le\zeta<1,
 \label{eq:wta-robust-centrality}
\end{equation}
put
$\widetilde p_G=\tr(\widetilde Z_1)/\sum_g\tr(\widetilde Z_g)$.
The denominator need not equal one under
\eqref{eq:wta-robust-centrality}.  Then
\eqref{eq:cond-r3}--\eqref{eq:cond-r5} specialize to
\begin{equation}
 \widetilde p_G\ge
 1-\frac{(1+\eta)(45/2)\alpha}{1-\zeta}.
 \label{eq:wta-robust-mass}
\end{equation}
The approximate transition lies between
$(1+\eta)(45/2)\alpha=1-\zeta$ and
$(1-\eta)(45/2)\alpha=1+\zeta$.
Equation \eqref{eq:cond-r13} also gives
$\kappa(\widetilde{\mathcal H}|_{\mathcal T})=\Omega(G)$ whenever
$\widetilde p_G$ is bounded below.  Trace error $\delta$ lowers the probability of observing the winner label
by at most $\delta$.

\subsection{Preparation, visibility, and extraction}
\label{sec:wta-central-state}

The root and physical central densities are
\begin{equation}
 \rho_\tau=\bigoplus_g Z_g(\tau),
 \qquad
 \rho_\tau^{\rm phys}=\frac1P\bigoplus_{g,i}
 R_{g,i}Z_g(\tau)R_{g,i}^\top.
 \label{eq:wta-central-densities}
\end{equation}
Both return the unique winner with probability $p_G$ when the group label is
measured.  The ordering
\eqref{eq:wta-resolvent-order} proves $p_G\ge1/G$, so the general upper
bound $\max\{p_G,1/G\}$ in \eqref{eq:cond-visibility} becomes $p_G$.  Thus
\cref{thm:cond-visibility} gives
\begin{align}
 p_G-\frac1G&\le\Dtr(\rho_0,\rho_1)\le p_G,
 \label{eq:wta-visibility}\\
 c_1-c_0&=\frac{(Gp_G-1)^2}{G(G-1)},\qquad c_0=\frac1G.
 \label{eq:wta-collision}
\end{align}

These formulas bound the cost of the first preparation through a collision
test on two independent preparations, with no query to verify a candidate
winner.  If $\tau\le1/(90G)$, \eqref{eq:wta-finite-mass} gives $p_G\ge3/4$,
and \eqref{eq:wta-collision} gives $c_1-c_0\ge1/8$ for every $G\ge2$.
Two independently prepared states of trace error $\epsilon$ change the
collision gap $c_1-c_0$ by at most $4\epsilon$.  Thus every fixed $\epsilon<1/32$ leaves
a constant gap.  At $\tau\le1/(135G)$, one has $p_G\ge5/6$; trace error
$1/20$ leaves gap at least $2/9-1/5=1/45$.  \Cref{thm:cond-unique-or} then implies
\begin{equation}
 \boxed{Q_\rho=\Omega(N\sqrt G)}
 \label{eq:wta-first-use-lower}
\end{equation}
for a first root or physical central-density preparation at either stated
schedule and tolerance.  At the first schedule, the physical duality gap is $3G\tau\le1/30$.  Under
\cref{def:amortized-accounting}, \eqref{eq:wta-first-use-lower} includes
setup queries; it does not cover an algorithm given a parity table or a
central-state oracle, which already contains the aggregated input.

Under the promise of exactly one winner, there is a matching coherent
upper bound.  A clean parity evaluator costs $T_{\rm eval}=N$ raw queries.
Instantiating
\eqref{eq:cond-preparation-upper}--\eqref{eq:cond-preparation-upper-p} in
\cref{prop:cond-rejection-preparation} gives
\begin{equation}
 Q_U=O\!\left(N(\sqrt{Gp_G}+1)\right),
 \label{eq:wta-preparation-upper}
\end{equation}
because the normalized $3$-dimensional winner and loser resolvents are public
functions of $(G,\tau)$.  The acceptance probability is classically known,
so exact phase-matched amplitude amplification removes any logarithmic
overhead from the raw-query count; the conditional local preparation cost
$T_{\rm loc}$ is gate cost only.  With the lower bound
\eqref{eq:cond-prep-tradeoff}, this is tight when the amplification
factor $\sqrt{Gp_G}$ diverges and the error resolves the winner mass $p_G$.  Hence
\begin{equation}
 Q_U=\begin{cases}
 \Theta(N\sqrt G),&0<\alpha<2/45,\\
 \Theta(NG^{1/4}),&\alpha=2/45,\\
 O(N),&\alpha>2/45.
 \end{cases}
 \label{eq:wta-preparation-phases}
\end{equation}
The extension to $k$ promised winners,
\eqref{eq:cond-ap3}--\eqref{eq:cond-ap6} in
\cref{sec:condensation-exact-k-preparation}, is not needed here.

Given an exact central-state preparation,
\eqref{eq:cond-algorithm-phases} gives these costs for extracting the
winner label:
\begin{center}
\begin{tabular}{@{}lccc@{}}
\toprule
regime & $p_G$ & verified copies & coherent calls \\ \midrule
$0<\alpha<2/45$ & $\Theta(1)$ & $\Theta(1)$ & $\Theta(1)$ \\
$\alpha=2/45$ & $\Theta(G^{-1/2})$ & $\Theta(G^{1/2})$ & $\Theta(G^{1/4})$ \\
$\alpha>2/45$ & $\Theta(G^{-1})$ & $\Theta(G)$ & $\Theta(G^{1/2})$ \\
\bottomrule
\end{tabular}
\end{center}
Verifying a label costs another $N$ raw queries.  Coherent
amplification needs a preparation unitary, its inverse, and a clean marker;
one trace-close output state does not suffice
(\cref{sec:state-conversion} treats that case).  The table counts uses of an
existing preparation, whereas \eqref{eq:wta-first-use-lower} bounds the
first preparation.

Visibility also shows when a fixed error allows a zero-query answer.  If
the allowed root trace error is at least $p_G$, then, under the promise of
zero or one winner, the public all-even root center is a valid output, by \eqref{eq:wta-resolvent-order} and
\eqref{eq:cond-visibility}.  At criticality this
happens for $G=\Omega(\epsilon^{-2})$, and for fixed supercritical $\alpha$
it happens for $G=\Omega_\alpha(\epsilon^{-1})$.  In physical coordinates,
lifting that state through the input's prefix products $R_{g,i}$ costs
exactly $N$ phase queries and preserves trace distance, so
\begin{equation}
 p_G\le\epsilon\quad\Longrightarrow\quad
 Q_\epsilon(\rho_{\rm central}^{\rm root})=0,
 \qquad Q_\epsilon(\rho_{\rm central}^{\rm phys})\le N.
 \label{eq:wta-zero-query-boundary}
\end{equation}
For sufficiently large fixed $\alpha$, even the public maximally mixed
physical state is zero-query: with $y_\alpha$ from
\eqref{eq:wta-even-root}, it suffices that
\begin{equation}
 \frac13-\frac{\alpha}{y_\alpha+3/10}+p_G\le\epsilon.
 \label{eq:wta-public-physical-boundary}
\end{equation}
No $\Omega(N)$ physical-state lower bound is claimed above the transition.

For algorithm design, condensation reveals the winner only in central
states whose preparation already costs as much as search, and a constant
winner mass forces $\kappa_{\rm red}$ linear in $G$.  The conservation law
\eqref{eq:cond-conservation} and the direct-search bounds
\cref{eq:cond-direct-search,eq:cond-search-crossover} therefore favor
search-then-crossover (find the winner, then solve only its group), at
cost
\begin{equation}
 O(N\sqrt G)+T_{\rm component\ solve}+O(N)_{\rm recovery}.
 \label{eq:wta-architecture}
\end{equation}
This family does not exhibit a QIPM speedup.

\subsection{Search then solve}
\label{sec:wta-search}

Write $\sigma_{g,j}=(-1)^{z_{g,j}}$.  With the XOR-oracle answer register in
$\ket-= (\ket0-\ket1)/\sqrt2$, query the public indices $j=1,\ldots,N$ in
order while retaining the group address.  Phase kickback implements
\begin{equation}
 \ket g\longmapsto\prod_{j=1}^N(-1)^{z_{g,j}}\ket g=h_g\ket g
 \label{eq:wta-clean-marker}
\end{equation}
with exactly $N$ raw queries and no parity or prefix garbage.

\begin{theorem}[Search-then-solve upper bound]
\label{thm:wta-search-upper}
For failure probability $\delta$, the value of the SDP can be returned using
$O(N\sqrt G\log(1/\delta))$ raw queries.  The same procedure returns a verified winning label, or `NONE' on an
all-even input, with a compact description of the corresponding exact root
optimizer; another $O(N)$ queries prepare that optimizer's physical state
instead.  It needs no QRAM.
\end{theorem}

\begin{proof}
Apply bounded-error quantum search to the clean marker
\eqref{eq:wta-clean-marker}, and verify a returned label by evaluating its
parity.  Output $14/5$ after successful verification and $29/10$ when the
search reports `NONE'.  Standard reversible addressing uses
$O(N\sqrt G\log(GP)\log(1/\delta))$ elementary routing gates and
$O(\log G+\log P)$ work qubits, apart from raw-query implementation.

If $g_*$ is odd, take $Z_{g_*}=r_-r_-^\top$ with
$r_-=(1,-1,1)^\top/\sqrt3$; if no odd group exists, take $g_*=1$ and
$Z_{g_*}=(I-r_+r_+^\top)/2$ with $r_+=(1,1,1)^\top/\sqrt3$.  Set every other
root to zero.  This is a compact optimizer output in the sense of
\cref{def:compact-optimizer-output}.  Streaming the $N$ prefix signs of
$g_*$ lifts it through \eqref{eq:wta-root-lift} and prepares the normalized
physical direct-sum state.  Listing its $P=33N$ nonzero blocks instead costs
$\Theta(P)$ output work.
\end{proof}

Together, \cref{thm:wta-search-upper} and
\eqref{eq:wta-quantum-lower} show that bounded-error value
estimation and winner recovery cost $\Theta(N\sqrt G)$, using standard search
\cite{boyer1998-tight-bounds-on-quantum-searching,
durr1996-a-quantum-algorithm-for-finding}.

Exact quantum query complexity depends on the promise.  Under
the promise of zero or one winner,
phase-matched exact amplification followed by exact parity verification gives
\begin{equation}
 Q_E(\operatorname{OR}^{0,1}_G\circ\operatorname{PARITY}_N)
 =\Theta(N\sqrt G).
 \label{eq:wta-exact-promised}
\end{equation}
Without that promise, the unique multilinear output polynomial has a
nonzero degree-$NG$ monomial, so the polynomial method gives an $NG/2$
lower bound; computing every parity gives
\begin{equation}
 Q_E(\operatorname{OR}_G\circ\operatorname{PARITY}_N)=\Theta(NG).
 \label{eq:wta-exact-unrestricted}
\end{equation}
If exactly one winner is promised, the decision is trivial, but finding
the winner's label still costs
$\Theta(N\sqrt G)$ for $G\ge2$.

\begin{remark}[Stronger access and setup costs]
\label{rem:wta-access-scope}
An oracle for $h_g$ removes the factor $N$ only because it is a stronger
input.  Parallel query ports may reduce depth but not the total number of
queries.  The accumulator chain is public.  The end-to-end bound requires
compact optimizer output; writing out $G$ arrays costs their output size.  If preprocessing uses $S$ raw queries and
later calls use $q_t$, any procedure that outputs the value or a decodable
winner obeys
\begin{equation}
 S+\sum_tq_t=\Omega(N\sqrt G).
 \label{eq:wta-setup-charge}
\end{equation}
In particular, the queries needed to build an input-dependent block
encoding, QRAM parity table, prefix table, null-space basis, or the cost
matrix obtained by eliminating the path variables count toward $S$.
\end{remark}

\subsection{Formulation tradeoff and novelty}
\label{sec:wta-scope}

A single trace row on the roots has $3G$ nonzeros, but its lift from root
directions to the matrix variables is an isometry.  The accumulator
chain \eqref{eq:wta-accumulator} reduces row sparsity to five and column
sparsity to two without changing the root SDP or the root-coordinate
Hessian, but its lift is not an isometry (\cref{sec:wta-public-start}).
Replacing the free scalars by differences of nonnegative variables, or
adding scalar log barriers, would change the geometry.  Thus we do not claim
a formulation with both bounded incidence and ambient condition number
one.

The query tools are prior work: adversary composition is due to
Høyer--Lee--Špalek, Belovs--Lee, and Reichardt; the randomized composition
framework includes Ben-David--Kothari; search and minimum finding are due to
BBHT and Dürr--Høyer; and fixed-point search is standard
\cite{hyer2005-tight-adversary-bounds-for-composite,
belovs2020-the-quantum-query-complexity-of,
reichardt2011-reflections-for-quantum-query-algorithms,
bendavid2018-randomized-query-complexity,
boyer1998-tight-bounds-on-quantum-searching,
durr1996-a-quantum-algorithm-for-finding,
yoder2014-fixed-point-quantum-search}.  Embeddings of Boolean functions into
conic optimal values, and sparsity-dependent coefficient-query lower bounds,
already appear in van
Apeldoorn et al., Theorem~29 and Corollary~30, and Apers--Gribling,
Theorem~8.4
\cite{apeldoorn2020-quantum-sdp-solvers,
apers2026-quantum-speedups-linear-programming-interior}.  Fawzi's non-SOC
theorem and the Rayleigh--Ritz trace identity are likewise established
tools.

To our knowledge, what is new is only their combination: a
bounded-incidence public accumulator chain for the trace normalization, an
exact spectral minimum that encodes
$\operatorname{OR}\circ\operatorname{PARITY}$ with a constant gap, a
public start whose root Hessian has condition number one, the exact
constants in the condensation transition, and the preparation costs in each
phase.  None of Grover search, parity, composition, the Rayleigh--Ritz trace
identity, generic SDP ill-conditioning, or log-det resolvents is claimed as
new by itself.
